\chapter{Introduction}

\section{Computer Network}
A computer network is a system that connects two or more computing devices for transmitting and sharing information. Computing devices include everything from a mobile phone to a server. These devices are connected using physical wires such as fiber optics, but they can also be wireless.

\subsection{Characteristics of a Computer Network}
\begin{enumerate}[itemsep=2ex]
   \item \textbf{Fault Detection}
   Fault detection ensures the reliability and integrity of data transmission by identifying and correcting errors in the network.

   \begin{itemize}
      \item \textbf{Data Link Layer:} Detects errors using CRC and retransmits damaged frames.
      \item \textbf{Transport Layer:} Uses TCP protocol for end-to-end error detection and retransmission.
   \end{itemize}
   
   \item \textbf{Scalability of Networks}
   Scalability ensures that a network can can handle growth in users, devices, and data traffic without affecting performance.
   \begin{itemize}
      \item \textbf{Data Link Layer:} Uses IP addressing and dynamic routing protocols to handle larger networks.
      \item \textbf{Transport Layer:} Handles multiple connections, congestion control (TCP/UDP).
   \end{itemize}
   
   \item \textbf{Quality of Service (QoS)}
   Quality of Service prioritizes network traffic to ensure smooth communication, reduce latency, and allocate bandwidth efficiently.
   \begin{itemize}
      \item \textbf{Network Layer:} Uses QoS protocols like DiffServ and MPLS for traffic prioritization.
      \item \textbf{Transport Layer:} Implements flow control and prioritization (e.g., TCP).
   \end{itemize}
   
   \item \textbf{Importance of Security}
   Security protects data confidentiality, integrity, and availability from unauthorized access and attacks.
   \begin{itemize}
      \item \textbf{Application Layer:} Enforces authentication and secure communication (HTTPS, VPNs).
      \item \textbf{Transport Layer:} Uses TCP protocol for end-to-end error detection and retransmission.
      \item \textbf{Network Layer:} Implements IPsec for secure packet transmission.
      \item \textbf{Session Layer:} Manages secure connections.
   \end{itemize}
\end{enumerate}


\section{OSI Layers}

\begin{tabularx}{\textwidth}{|l|l|p{3cm}|X|}\hline
   & \thc{Data Unit} & \thc{Layer} & \thc{Function}\\\hline
   
   \rspan{4}{Host Layers} & \rspan{3}{Data}
   & 7. Application\newline(Desktop Layer)&
   Network Virtual Terminal, File Transfer Access and Management (FTAM), Mail Services, Directory Services
   \\\cline{3-4}

   & & 6. Presentation &
   Data Translation, Encryption/Decryption, Compression
   \\\cline{3-4}

   & & 5. Session &
   Session Establishment and Termination, Dialog Control, Synchronization
   \\\cline{2-4}

   & Segments & 4. Transport\newline(Heart of OSI) &
   Segmentation and Reassembly, Error Detection and Correction, Flow Control, Port Addressing
   \\\hline

   \rspan{3}{Media Layers}
   & Packet & 3. Network &
   Routing, Logical Addressing
   \\\cline{2-4}

   & Frame & 2. Data Link & Framing, Physical Addressing, Error Control, Flow Control, Access Control
   \\\cline{2-4}

   & Bit & 1. Physical &
   Bit Synchronization, Bit Rate Control, Physical Topology, Transmission Modes
   \\\hline
\end{tabularx}

